% Options for packages loaded elsewhere
\PassOptionsToPackage{unicode}{hyperref}
\PassOptionsToPackage{hyphens}{url}
\PassOptionsToPackage{dvipsnames,svgnames,x11names}{xcolor}
\documentclass[
  10pt,
]{article}
\usepackage{xcolor}
\usepackage[margin=2.4cm]{geometry}
\usepackage{amsmath,amssymb}
\setcounter{secnumdepth}{-\maxdimen} % remove section numbering
\usepackage{iftex}
\ifPDFTeX
  \usepackage[T1]{fontenc}
  \usepackage[utf8]{inputenc}
  \usepackage{textcomp} % provide euro and other symbols
\else % if luatex or xetex
  \usepackage{unicode-math} % this also loads fontspec
  \defaultfontfeatures{Scale=MatchLowercase}
  \defaultfontfeatures[\rmfamily]{Ligatures=TeX,Scale=1}
\fi
\usepackage{lmodern}
\ifPDFTeX\else
  % xetex/luatex font selection
  \setmainfont[]{Libertinus Serif}
  \setmonofont[Scale=0.85]{Latin Modern Mono}
  \setmathfont[]{Latin Modern Math}
\fi
% Use upquote if available, for straight quotes in verbatim environments
\IfFileExists{upquote.sty}{\usepackage{upquote}}{}
\IfFileExists{microtype.sty}{% use microtype if available
  \usepackage[]{microtype}
  \UseMicrotypeSet[protrusion]{basicmath} % disable protrusion for tt fonts
}{}
\makeatletter
\@ifundefined{KOMAClassName}{% if non-KOMA class
  \IfFileExists{parskip.sty}{%
    \usepackage{parskip}
  }{% else
    \setlength{\parindent}{0pt}
    \setlength{\parskip}{6pt plus 2pt minus 1pt}}
}{% if KOMA class
  \KOMAoptions{parskip=half}}
\makeatother
\usepackage{longtable,booktabs,array}
\newcounter{none} % for unnumbered tables
\usepackage{calc} % for calculating minipage widths
% Correct order of tables after \paragraph or \subparagraph
\usepackage{etoolbox}
\makeatletter
\patchcmd\longtable{\par}{\if@noskipsec\mbox{}\fi\par}{}{}
\makeatother
% Allow footnotes in longtable head/foot
\IfFileExists{footnotehyper.sty}{\usepackage{footnotehyper}}{\usepackage{footnote}}
\makesavenoteenv{longtable}
\setlength{\emergencystretch}{3em} % prevent overfull lines
\providecommand{\tightlist}{%
  \setlength{\itemsep}{0pt}\setlength{\parskip}{0pt}}

\providecommand{\E}{\mathbb{E}}
\providecommand{\N}{\mathbb{N}}
\providecommand{\Z}{\mathbb{Z}}
\providecommand{\R}{\mathbb{R}}
\providecommand{\eps}{\varepsilon}
\providecommand{\ceil}[1]{\left\lceil #1\right\rceil}
\providecommand{\floor}[1]{\left\lfloor #1\right\rfloor}
\AtBeginDocument{\let\setminus\smallsetminus}
\usepackage[most]{tcolorbox}
\newtcolorbox{warning}{colback=red!5,colframe=red!65!black,boxrule=0.8pt,arc=2pt,left=6pt,right=6pt,top=5pt,bottom=5pt}
\usepackage{etoolbox}
\AtBeginEnvironment{longtable}{\small}
\setlength{\emergencystretch}{3em}
\usepackage{fancyhdr}
\pagestyle{fancy}
\fancyhf{}
\fancyhead[L]{\small\bfseries UNREFEREED GPT-6 ASTRA OUTPUT}
\fancyhead[R]{\small\thepage}
\renewcommand{\headrulewidth}{0.4pt}
\usepackage{bookmark}
\IfFileExists{xurl.sty}{\usepackage{xurl}}{} % add URL line breaks if available
\urlstyle{same}
\hypersetup{
  pdftitle={Informal Question (Section 2)},
  pdfauthor={Writeup: gpt-6-astra (single model pass)},
  colorlinks=true,
  linkcolor={blue!50!black},
  filecolor={Maroon},
  citecolor={Blue},
  urlcolor={blue!50!black},
  pdfcreator={LaTeX via pandoc}}

\title{Informal Question (Section 2)}
\usepackage{etoolbox}
\makeatletter
\providecommand{\subtitle}[1]{% add subtitle to \maketitle
  \apptocmd{\@title}{\par {\large #1 \par}}{}{}
}
\makeatother
\subtitle{Unrefereed candidate disproof by GPT-6 Astra}
\author{Writeup: \texttt{gpt-6-astra} (single model pass)}
\date{Catalog id \texttt{1812.09215\_\_00} --- generated 2026-09-17}

\begin{document}
\maketitle

\begin{warning}

\textbf{UNREFEREED MODEL OUTPUT.} This document was selected solely
because GPT-6 Astra labelled its own result
\texttt{would\_publish:\ true} and returned \texttt{proved} or
\texttt{disproved}. It has not passed the adversarial LLM referee used
for the earlier Sol campaign, has not been checked by a human
mathematician, and has not been checked for novelty. Treat every
mathematical and bibliographic claim below as unverified.

\end{warning}

\section{Summary and provenance}\label{summary-and-provenance}

{\def\LTcaptype{none} % do not increment counter
\begin{longtable}[]{@{}
  >{\raggedright\arraybackslash}p{(\linewidth - 2\tabcolsep) * \real{0.5000}}
  >{\raggedright\arraybackslash}p{(\linewidth - 2\tabcolsep) * \real{0.5000}}@{}}
\toprule\noalign{}
\endhead
\bottomrule\noalign{}
\endlastfoot
Catalog id & \texttt{1812.09215\_\_00} \\
Catalog entry &
\href{https://graph-theory-ai.github.io/graph-conjectures/arxiv/1812.09215__00/}{Informal
Question (Section 2)} \\
Source corpus & arxiv \\
Campaign leg & selected second attempts (\texttt{attacks\_retry}) \\
Source paper / entry & Lipschitz bijections between boolean functions \\
Model verdict & \textbf{disproved} (confidence: high) \\
Model's one-line claim & No: every such bijection satisfies n \textless=
1+(k-1)+\ldots+(k-1)\^{}(D-1), where D is its inverse Lipschitz
constant, and linear 2-Lipschitz tree maps attain the resulting
bound. \\
Model & \texttt{gpt-6-astra}, reasoning effort \texttt{max},
\texttt{mode=pro}, flex service tier \\
Original artifact &
\texttt{attacks\_retry/1812.09215\_\_00/output.md} \\
Independent review & \textbf{None} \\
Model's caveats & Uses full-cube bijections and global k-junta locality;
literature priority has not been independently checked. \\
\end{longtable}
}

\section{Problem statement}\label{problem-statement}

\emph{No extracted statement was stored in the artifact metadata.}

\section{Astra writeup}\label{astra-writeup}

\emph{The text below is the model output verbatim, apart from moving its
machine-readable verdict block into the summary above and shifting
Markdown heading levels for this document. Mechanical TeX defects and
equation tags were normalized where needed for compilation.}

\subsection{A sharp negative answer}\label{a-sharp-negative-answer}

The gap closes on the lower-bound side. In fact, the forward Lipschitz
assumption is unnecessary for the lower bound, and a rooted-tree
construction gives the exact optimum.

The proof below does not rely on the previous attempt's
localized-dependency estimate. It retains the elementary inverse-edge
setup, verified below, but replaces dependency-graph packing by a
bounded-branching count of possible target outputs.

\subsubsection{1. Statement}\label{statement}

Let \(Q_n=\{0,1\}^n\), equipped with Hamming distance. A Dictator-to-XOR
bijection is a bijection \[
\phi:Q_n\longrightarrow Q_n
\] such that \[
\bigoplus_{a=1}^n\phi_a(x)=x_1
\qquad\text{for every }x\in Q_n.                     \qquad\text{(1)}
\]

Assume that, for each output coordinate \(a\), there is a set
\(P_a\subseteq[n]\), with \(|P_a|\le k\), such that \(\phi_a\) depends
only on the input coordinates in \(P_a\). Write \[
D=\operatorname{Lip}(\phi^{-1}).
\] On a Hamming cube, the Lipschitz constant equals the maximum stretch
of an edge, so \(D\) is an integer.

For \(k\ge2\) and \(d\ge1\), put \[
N_k(d)=\sum_{r=0}^{d-1}(k-1)^r.
\]

\paragraph{Theorem}\label{theorem}

Every bijection satisfying (1) and the locality condition obeys \[
\boxed{\quad n\le N_k(D).\quad}                       \qquad\text{(2)}
\]

Conversely, whenever \(n\le N_k(d)\), there is a \textbf{linear,
2-Lipschitz} Dictator-to-XOR bijection whose output coordinates depend
on at most \(k\) inputs and whose inverse is \(d\)-Lipschitz.

Consequently, the minimum possible inverse Lipschitz constant is exactly
\[
\boxed{
D_{\min}(n,k)=
\begin{cases}
n, & k=2,\\[2mm]
\displaystyle
\left\lceil
\frac{\log((k-2)n+1)}{\log(k-1)}
\right\rceil, & k\ge3.
\end{cases}}                                         \qquad\text{(3)}
\] This remains the exact optimum if any fixed forward Lipschitz bound
\(C\ge2\) is imposed.

In particular, \[
\boxed{\quad D\ge 1+\frac{\log n}{\log k}.\quad}       \qquad\text{(4)}
\] Thus the inverse cannot be \(o(\log n/\log k)\)-Lipschitz.

\subsubsection{2. The counting argument}\label{the-counting-argument}

We prove a slightly stronger pointwise statement.

Fix an input \(x\in Q_n\), and write \(y=\phi(x)\). For each output
coordinate \(a\), define \[
x^{(a)}=\phi^{-1}(y\oplus e_a),
\qquad
R_a=\{i:x_i\ne x_i^{(a)}\}.
\] The sets \(R_a\) are nonempty and pairwise distinct.

Equation (1) gives \[
x^{(a)}_1
=\bigoplus_b (y\oplus e_a)_b
=\left(\bigoplus_b y_b\right)\oplus1
=x_1\oplus1.
\] Hence \[
1\in R_a\qquad\text{for every }a.                    \qquad\text{(5)}
\]

For a set \(T\subseteq[n]\), let \[
x^T=x\oplus\bigoplus_{i\in T}e_i.
\] Thus \(x^{R_a}=x^{(a)}\).

Fix an integer \(d\ge1\). For nonempty \(T\subseteq[n]\), with
\(|T|\le d\), define the family of target labels \[
A_d(T)=\{a:T\subseteq R_a,\ |R_a|\le d\}.
\]

\paragraph{Claim}\label{claim}

For every such \(T\), \[
\boxed{\quad
|A_d(T)|\le N_k(d-|T|+1).
\quad}                                                \qquad\text{(6)}
\]

\paragraph{Proof}\label{proof}

We use induction on \(d-|T|\).

If \(|T|=d\), every \(a\in A_d(T)\) must satisfy \(R_a=T\). Since the
\(R_a\) are distinct, \[
|A_d(T)|\le1=N_k(1).
\]

Now suppose \(|T|<d\). Because \(T\ne\varnothing\) and \(\phi\) is
injective, \[
\phi(x^T)\ne\phi(x).
\] Choose an output coordinate \(b\) on which these two images differ:
\[
\phi_b(x^T)\ne\phi_b(x).                              \qquad\text{(7)}
\] Since \(\phi_b\) depends only on \(P_b\), equation (7) implies \[
P_b\cap T\ne\varnothing,
\qquad\text{and therefore}\qquad
|P_b\setminus T|\le k-1.                              \qquad\text{(8)}
\]

Consider any target label \(a\in A_d(T)\) with \(a\ne b\). By the
definition of \(R_a\), \[
\phi(x^{R_a})=\phi(x)\oplus e_a.
\] Since \(a\ne b\), it follows that \[
\phi_b(x^{R_a})=\phi_b(x)\ne\phi_b(x^T).              \qquad\text{(9)}
\] The inputs \(x^{R_a}\) and \(x^T\) differ precisely on
\(R_a\setminus T\), because \(T\subseteq R_a\). Locality and (9)
therefore imply \[
P_b\cap(R_a\setminus T)\ne\varnothing.
\] Thus \(a\) belongs to \(A_d(T\cup\{i\})\) for at least one
\(i\in P_b\setminus T\).

We have proved the covering relation \[
A_d(T)
\subseteq
\{b\}\ \cup\
\bigcup_{i\in P_b\setminus T}A_d(T\cup\{i\}).         \qquad\text{(10)}
\] The exceptional set \(\{b\}\) has size one; it does not matter
whether \(b\) actually belongs to \(A_d(T)\).

Applying the induction hypothesis to every child set \(T\cup\{i\}\), and
using (8), gives \begin{equation*}
\begin{aligned}
|A_d(T)|
&\le 1+\sum_{i\in P_b\setminus T}|A_d(T\cup\{i\})|\\
&\le 1+(k-1)N_k(d-|T|)\\
&=N_k(d-|T|+1).
\end{aligned}
\end{equation*} This completes the induction. \(\square\)

\paragraph{Applying the claim}\label{applying-the-claim}

By (5), every inverse edge at \(y\) changes input coordinate \(1\).
Therefore \[
A_d(\{1\})=\{a:|R_a|\le d\}.
\] Taking \(T=\{1\}\) in (6), we obtain the pointwise bound \[
\boxed{\quad
\#\left\{a:
d_H\bigl(\phi^{-1}(y),\phi^{-1}(y\oplus e_a)\bigr)
\le d
\right\}
\le N_k(d).
\quad}                                                \qquad\text{(11)}
\]

For \(d=D\), all \(n\) inverse edges have stretch at most \(D\).
Equation (11) yields \[
n\le N_k(D),
\] proving (2).

The essential distinction from the earlier packing argument is this: at
a partial input change \(T\), choose \textbf{one} currently changed
output \(b\). Its own target label accounts for at most one possibility.
Every other target must add an input from the set \(P_b\setminus T\),
which has size at most \(k-1\). This gives a branching factor \(k-1\),
without any dependence on the radius \(d\).

\subsubsection{3. Sharpness: the rooted-tree
construction}\label{sharpness-the-rooted-tree-construction}

Let \(T\) be a rooted tree with \(n\) vertices, root \(1\), maximum
number of children at most \(k-1\), and height at most \(d-1\).

Such a tree exists whenever \[
n\le 1+(k-1)+\cdots+(k-1)^{d-1}:
\] take the first \(n\) vertices in breadth-first order in the complete
rooted \((k-1)\)-ary tree of height \(d-1\).

Index both input and output coordinates by the vertices of \(T\). Define
\[
\phi_v(x)
=
x_v\oplus\bigoplus_{w:\,w\text{ is a child of }v}x_w. \qquad\text{(12)}
\]

We verify all the required properties.

\paragraph{Locality}\label{locality}

Each \(\phi_v\) depends on its own input and on its children's inputs,
hence on at most \[
1+(k-1)=k
\] input coordinates.

\paragraph{Dictator-to-XOR condition}\label{dictator-to-xor-condition}

In the XOR of all outputs, the root input appears once. Every nonroot
input appears twice: in its own output and in its parent's output.
Consequently, \[
\bigoplus_v\phi_v(x)=x_1.
\]

\paragraph{Bijectivity}\label{bijectivity}

The inverse is \[
x_v
=
\bigoplus_{w\in\operatorname{subtree}(v)}y_w,
\qquad y=\phi(x).                                    \qquad\text{(13)}
\] Indeed, summing (12) over the subtree rooted at \(v\) cancels every
input in that subtree except \(x_v\). Thus every input is recovered
uniquely.

\paragraph{Forward Lipschitz constant}\label{forward-lipschitz-constant}

Flipping the root input changes only the root output. Flipping a nonroot
input \(x_v\) changes exactly the outputs indexed by \(v\) and its
parent. Hence \[
\operatorname{Lip}(\phi)\le2.
\]

\paragraph{Inverse Lipschitz constant}\label{inverse-lipschitz-constant}

By (13), flipping output \(y_w\) changes exactly those inputs indexed by
ancestors of \(w\), including \(w\). Its inverse stretch is therefore \[
\operatorname{depth}(w)+1\le d.
\] Thus \[
\operatorname{Lip}(\phi^{-1})\le d.
\]

This proves the converse assertion. Choosing the smallest \(d\) with
\(n\le N_k(d)\) proves the exact optimum (3).

For completeness, when \(n\ge2\), a forward Lipschitz bound below \(2\)
is impossible: flipping an input coordinate other than \(1\) preserves
output parity, so injectivity forces the output Hamming distance to be a
positive even integer. Therefore the construction also uses the smallest
possible forward Lipschitz constant.

\subsubsection{4. Consequences for the
question}\label{consequences-for-the-question}

For \(k=2\), equation (2) is simply \[
n\le D.
\]

For \(k\ge3\), \[
n\le\frac{(k-1)^D-1}{k-2},
\] which gives the second line of (3).

A convenient uniform consequence follows from \[
\sum_{r=0}^{D-1}(k-1)^r
\le
\sum_{r=0}^{D-1}\binom{D-1}{r}(k-1)^r
=k^{D-1}.
\] Thus \[
n\le k^{D-1},
\qquad\text{so}\qquad
D\ge1+\frac{\log n}{\log k}.
\] This excludes \[
D=o\!\left(\frac{\log n}{\log k}\right)
\] along every sequence with \(n,k\to\infty\).

The degenerate locality case \(k=1\) permits no such bijection for
\(n\ge2\): in (10), \(P_b\setminus T\) is empty, so taking \(T=\{1\}\)
gives \(n\le1\).

\textbf{Conclusion.} Under the standard full-cube and global-junta
definitions in the question, nonlinear maps offer no improvement in the
optimal maximum inverse stretch: linear 2-Lipschitz rooted-tree maps are
already exactly optimal. No additional conjecture or unproved
mathematical step is used. The remaining caveat is literature priority,
which has not been independently checked here.

\end{document}
